\section{Two sign-exchanged kernel types with repeated multiplicity}
\label{sec:heisenberg-two-isotype-symmetric}

\paragraph{Author freeze and origin.}
D-TD39, reserved question Q2eu, is a derivation by
\texttt{/root/imprimitive} from a lead of \texttt{/root/sixfolds}
and the whole symmetric-quotient budget proposed by \texttt{/root}.
The two-isotype irreducibility, faithful multiplicity actions,
joint kernel bound and actual normal subgroup construction are
proved below.  D-TD37 supplies a separately stated source
obstruction once that actual subgroup has been constructed.
Complete author proof and source-interface readings pass.
Independent and root whole-proof reviews of this section also pass.
In particular no sphere-cover hypothesis is transferred to the
index-two subgroup.

Use the actual stable whole Heisenberg factor setup of D-TD05:
\[
 \begin{gathered}
 C\longrightarrow\mathbf P^1\text{ is a connected }J\text{-Galois cover},
 \quad V=W\otimes R\text{ is a faithful finite complex coefficient space},
 \quad e=\dim V,\\
 P=\operatorname{im}(J\to\operatorname{PGL}(W)),\qquad
 Q\triangleleft P,\quad Q\simeq\mathcal Q\simeq\mathbf F_p^{2m},
 \quad P/Q=H\le\operatorname{Sp}(\mathcal Q),\quad \dim W=p^m.
 \end{gathered}
 \tag{TSI.1}
\]
The projective Q-action is irreducible with nondegenerate
commutator form.  The defining H-action on \(\mathcal Q\) is
faithful and irreducible over the original \(\mathbf F_p\).
R is the entire external partner, with no dimension or support
condition.  Assume the actual label image has a surjection
\[
 \psi:H\twoheadrightarrow S_n,\qquad n\ge6,
 \qquad N=\ker\psi,\qquad H_0=\psi^{-1}(A_n),
 \tag{TSI.2}
\]
and that the restriction of \(\mathcal Q\) to N has exactly
two distinct original-field isotypic components.

\begin{theorem}[D-TD39]
In (TSI.1)--(TSI.2), the two components have the form
\[
 \mathcal Q=\mathcal Q_1\oplus\mathcal Q_2,
 \qquad \mathcal Q_i\big|_N\simeq S_i^{\oplus r},
 \qquad S_1\not\simeq S_2,
 \qquad \dim_{\mathbf F_p}S_1=\dim_{\mathbf F_p}S_2=d_0.
 \tag{TSI.3}
\]
Suppose \(r>1\).  If the actual original coefficient inertia
has moving-rank sum \(2e+6\) and \(e>504\), then
\[
                  p^{d_0}<28900,\qquad d_0\le14,
                  \qquad |N|<B:=28900^{28}.
 \tag{TSI.4}
\]
If \(n\ge B+2\), no such actual sphere source with that budget
exists, without needing to assume \(e>504\) separately.  More
precisely, from the presumed budget in that range one constructs
an actual subgroup
\[
 E\triangleleft H,\qquad E=E',\qquad
 \psi(E)=A_n,\qquad E\cap N=Z(E),\qquad E/Z(E)\simeq A_n,
 \tag{TSI.5}
\]
and it is this subgroup, not an abstract section, that supplies
the source contradiction through D-TD37.
\end{theorem}

\begin{proof}
\paragraph{Step 1: the two actual subspaces and their irreducibility.}
The N-socle of \(\mathcal Q\) is nonzero and H-invariant.
Original-field irreducibility therefore makes the restriction
semisimple, so the two canonical isotypic spaces are genuine
direct summands.  H permutes them transitively: otherwise each
would be H-invariant.  N fixes both, because its action preserves
every N-isotypic space.  Thus the component action factors through
\(S_n\) and is a nontrivial map to \(S_2\).
Its kernel is \(A_n\), since the abelianization of \(S_n\)
is its sign quotient.  Therefore every element of H0 fixes
each \(\mathcal Q_i\), and every element outside H0 exchanges
them.  Such an exchange preserves the dimensions and twists one
simple N-type into the other, establishing the common d0 and r
in (TSI.3).

Each \(\mathcal Q_i\) is irreducible for H0 over
\(\mathbf F_p\).  Indeed, suppose a nonzero proper subspace
\(U\subset\mathcal Q_1\) is H0-invariant, and choose
\(t\in H\setminus H_0\).  The subspace \(U\oplus tU\)
is H0-invariant because H0 is normal.  It is also t-invariant:
\(t^2\in H_0\), so \(t^2U=U\).  Since H is generated by
H0 and t, it would be a nonzero proper H-invariant subspace
of \(\mathcal Q\), a contradiction.
This argument does not assert that either summand is
nondegenerate for the symplectic form.  It also covers a dual
pair of totally isotropic components.

\paragraph{Step 2: the multiplicity quotient is actual and faithful.}
Fix i.  The image of \(\mathbf F_p[N]\) on the simple
\(S_i\) is one full matrix algebra over its finite commuting
field.  As in the proved algebra argument of D-TD27, write
\[
 F_i=\mathbf F_{p^{f_i}},\qquad d_0=a_i f_i,\qquad
 \mathbf F_p[N]\big|_{\mathcal Q_i}
      =\operatorname{Mat}_{a_i}(F_i)\otimes1_r,
 \qquad \mathcal Q_i\simeq F_i^{a_i}\otimes_{F_i}F_i^r.
 \tag{TSI.6}
\]
Here the algebra acts faithfully by its definition as an image;
N itself need not act faithfully on \(S_i\).
The finite-algebra proof uses a faithful simple image algebra,
not faithfulness of the abstract group N.

H0 normalizes this algebra and its center.  Its action on the
field \(F_i\) kills N and hence factors through the perfect
\(H_0/N=A_n\); a finite field has cyclic automorphism group,
so that action is trivial.  Thus H0 is \(F_i\)-linear.
Skolem--Noether and the exact commuting matrix algebra give
\[
 h\big|_{\mathcal Q_i}=A_i(h)\otimes B_i(h),\qquad
 \beta_i:A_n\longrightarrow\operatorname{PGL}_r(F_i),
       \qquad \beta_i(hN)=[B_i(h)].
 \tag{TSI.7}
\]
The tensor factors are unique up to reciprocal scalars, so the
factor map from H0 is a genuine projective homomorphism.  It
kills N and induces the displayed genuine homomorphism \(\beta_i\).
No quotient is inferred from a nonfaithful restriction image.

If \(\beta_i\) were trivial, every \(F_i\)-subspace
\(F_i^{a_i}\otimes L\), with L a multiplicity line, would
be H0-invariant.  Because \(r>1\), this is a nonzero proper
\(\mathbf F_p\)-subspace of \(\mathcal Q_i\), contrary
to Step 1.  The map is nontrivial and hence faithful by
simplicity of \(A_n\).  In particular, if an element of H0
acts trivially on \(\mathcal Q_i\), its natural image is
one, so its kernel lies in N.  The actual restriction image
therefore does inherit the A_n quotient, but only after this
argument establishes the needed kernel inclusion.

D-TD32 applied to the faithful projective map also gives
\[
                             r>n/5.
 \tag{TSI.8}
\]
All these field and matrix-algebra arguments are valid for an
arbitrary irreducible original-field module; no symplectic form
on \(\mathcal Q_i\) is used.

\paragraph{Step 3: every nonidentity natural image costs label rank.}
For \(h\in H_0\) with \(\psi(h)\ne1\), each matrix
\(B_i(h)\) is nonscalar.  Over an algebraic closure its
geometric projective rank is therefore at least one, including
when it is nonsemisimple.  D-TD27's geometric tensor inequality
and multiplication by the field degree \(f_i\) give
\[
 \operatorname{rank}_{\mathbf F_p}(1-h\mid\mathcal Q_i)
                       \ge a_i f_i=d_0.
 \tag{TSI.9}
\]
The two ranks add on \(\mathcal Q\).  For \(h\notin H_0\),
the operator exchanges the two equal-dimensional components.
A fixed vector is determined by its first component, so its
fixed space has dimension at most \(\dim\mathcal Q_1=r d_0\).
Its moving rank is consequently at least \(r d_0\ge d_0\).
The following conservative common bound is sufficient:
\[
       \psi(h)\ne1\quad\Longrightarrow\quad
       \operatorname{rank}_{\mathbf F_p}(1-h)\ge d_0.
 \tag{TSI.10}
\]
This includes both even and odd natural images.  It holds for
every original power according to the image of that power.

\paragraph{Step 4: the whole symmetric quotient, not an H0 base.}
Compose the actual maps \(J\to P\to H\to S_n\).
The corresponding connected S_n quotient cover of the sphere
has genus at least two.  Indeed \(n\ge6\), so S_n is
non-solvable of order greater than sixty; the accepted genus-zero
and genus-one exclusions in D-TD02 apply.  Its nonidentity
inertia images of orders \(o_i\) therefore satisfy
\(\sum_i(1-1/o_i)\ge85/42\).

For every full coefficient element whose S_n-image is nonidentity,
the exact Heisenberg trace identity and (TSI.10), with the whole
external partner retained, give normalized absolute trace at
most \(c=p^{-d_0/2}\).  As in D-TD28, take the full original
order L for each coefficient inertia.  If its natural image has
order o, exactly \(L/o\) powers have identity image and are
bounded by one.  All remaining powers are bounded by c.
Fourier inversion consequently gives every-target support at
least \((1-c)(1-1/o)\).  Summing and using the actual budget
gives
\[
          2+\frac6e\ge(1-c)\frac{85}{42},\qquad
          e(1-85c)\le252.
 \tag{TSI.11}
\]
No sphere-cover assertion has been made for H0 or for a
restriction to a component.  The quotient used here is the full
actual symmetric quotient of the original sphere cover.

If \(e>504\) and \(p^{d_0}\ge170^2=28900\), then
\(c\le1/170\) and (TSI.11) gives \(e\le504\), a
contradiction.  Thus \(p^{d_0}<28900<2^{15}\), which also
gives \(d_0\le14\).

Let \(N_i\) be the image of N on \(S_i\).  The kernels
of those two actions intersect trivially, because the full
N-action on the faithful \(\mathcal Q\) is the sum of
copies of those two actions.  Hence the joint map, though
neither individual map need be faithful, is injective:
\[
 N\lhook\joinrel\longrightarrow N_1\times N_2,
 \qquad
 |N|\le|N_1||N_2|
       <p^{2d_0^2}<28900^{2d_0}\le28900^{28}=B.
 \tag{TSI.12}
\]
This proves the quantitative assertions (TSI.4).

\paragraph{Step 5: construction normal in all of H and the source conclusion.}
Suppose \(n\ge B+2\).  By (TSI.8),
\(e\ge\dim W=p^{r d_0}\ge2^r>2^{n/5}>504\):
indeed \(B>10^6>50\), and \(2^{10}=1024>504\).
Thus Step 4 applies without an extra coefficient-degree premise.

Conjugation on N induces \(H_0/N=A_n\to\operatorname{Out}(N)\).
The elementary bound \(|\operatorname{Out}(N)|\le|N|!\le B!\)
and \(|A_n|=n!/2>B!\) exclude an injection, so simplicity
makes this map trivial.  Every H0-conjugation on N is inner.
Put
\[
 C_0=C_{H_0}(N),\qquad H_0=NC_0,\qquad
 C_0\cap N=Z(N),\qquad C_0/Z(N)\simeq A_n,
 \qquad E=C_0'.
 \tag{TSI.13}
\]
The group \(C_0=H_0\cap C_H(N)\) is normal in all of H,
since H0 and N are normal.  Therefore E is also normal in H.
Perfectness of \(A_n\) makes \(C_0=E Z(N)\).  This
Z(N) is central in C0, so
\(E=C_0'=[E Z(N),E Z(N)]=E'\).  Its image under \(\psi\)
is \(A_n\), and its kernel is central.  More precisely,
\(E\cap N\subseteq Z(N)\) because E centralizes N,
while a central element of E maps to the trivial center of
\(A_n\).  These inclusions give \(E\cap N=Z(E)\).
This proves every assertion of (TSI.5).

The actual matrix action is faithful on E; the full H-action is
still irreducible over \(\mathbf F_p\); and the original
\(\psi:H\to S_n\) is compatible with this actual normal
quasisimple subgroup.  Since \(n\ge B+2>10^6\), all
hypotheses of D-TD37 hold for the same entire Heisenberg factor
and external partner.  Its actual-source conclusion contradicts
the presumed moving-rank sum \(2e+6\).
\end{proof}

\paragraph{Retained boundary.}
The restriction has two distinct original-field N-types and each
occurs with multiplicity greater than one.  Two types of
multiplicity one remain outside this theorem.  The homogeneous
symmetric case is different: a multiplicity-two space can have
trivial alternating action and become irreducible only through
the odd operator.  Nothing here discards that possibility.
The theorem makes no claim about original complex character fields,
polarizations, source existence, or a symplectic form on either
individual N-isotypic space.
